\scanpage{25}
tels que $\forall j\geq1$, on ait comm. dans
\diagram{\begin{tikzcd}[ampersand replacement=\&,column sep=large,row sep=large]
M.T^{(i)} \arrow[r,"{.T^{(j)}}"] \arrow[d,"{u_i\otimes T^{(\unsure{i|h})}}"'] \& M T^{(i+j)} \arrow[d,"{u_{i+j}\otimes T^{(\unsure{i|h}+j)}}"]\\
M'.T^{(i)} \arrow[r,"{.T^{(\unsure{j|i})}}"] \& M'.T^{(i+j)}
\end{tikzcd}}
i.e. tels qu'on ait
\[
c_{ij}(u_{i+j}-u_i)=0\qquad\text{si}\quad i\geq N,\ j\geq1,
\]
et le th. nous permet de conclure que tous les $u_i$ sont égaux, d'où la conclusion
\[
\operatorname{Hom}_k(M,M')\xrightarrow{\sim}\operatorname{Hom}_{S\text{-mod gr.}\unsure{d|s|{}}}(\unsure{w|\omega}_N(M),w_N(M'))
\]

\uline{Rmq 1} Le foncteur $w_N$ admet un adjoint à droite $\pi_N$
\[
\operatorname{Hom}_{S\text{-mod gr.}d}(w_N(M),\deleted{\unsure{M'}}\mathcal M)\unsure{\simeq|\xrightarrow{\sim}}\operatorname{Hom}_k(M,\pi_N(\mathcal M))
\]
en posant
\[
\begin{aligned}
\pi_N(\mathcal M)&\subset\prod_{i\geq N}\mathcal M_i\\
\shortparallel&\\
\bigl\{(x_i)_{i\geq N}\ &\big|\ x_iT^{(j)}=c_{ij}x_{j+j}\quad\forall i\geq N,\ j\geq1\bigr\}
\end{aligned}
\]
En considération de ceci, \deleted{il existe} les parties a) et b) du corollaire 1 sont équivalentes, b) signifiant que l'hom can. d'adjonction
\[
M\longrightarrow\pi_N(w_N(M))
\]
est \deleted{\unsure{bij}} un isomorphisme.

\uline{Rmq 2} Considérons $w_0(M)=M\otimes_k S=M\otimes_k k\{T\}$, alors $w_N(M)=S^{+(N)}.w_0(M)$. \deleted{Donc \unsure{pour} en un sens, $w_0$ est} \inserted{On peut aussi écrire, si $N\geq N'\geq0$}
\[
w_N(M)\simeq\deleted{\unsure{\prod}}\sum_{i\geq N}w_{N'}(M)_i\overset{\mathrm{def}}{=}\underset{\text{(``troncature'')}}{\tau_{\geq N}}(w_{N'}(M)).
\]
Donc en
% archive sheet number: 24

